\section{Transporting certificates to original executions}\label{sec:transfer}

The direct replay arguments cover every positive width, but they do not
subsume two other conclusions: expansion from \emph{every} large output set
to the original crossing set, and storage capacity on an original run's entry
quotient. This appendix retains those conclusions. Their finite interfaces
use the actual aligned simulator of~\cite{HighRateStorage2026}; none identifies
its dependency graph with the original graph.

\subsection{A guarded projection}

Fix a halting source execution with $t\ge1$ transitions and an integer $s\ge1$.
The aligned simulator has $k$ corresponding work heads and halts after
\[
 B_0=10s+1,\qquad U=B_0(\lfloor t/s\rfloor+1)+s
\]
transitions. The extra complete terminal batch is present even when $s\mid t$.
Its final $s$ steps form a stationary drain. Define
\[
 q(u)=\min\bigl((\lfloor u/B_0\rfloor+1)s,t-1\bigr),\qquad 0\le u<U.
\]
Write $x_h$ for source positions and $y_h$ for physical simulator positions.
At batch $j$ the common anchor is
$s\lfloor x_h(js)/s\rfloor$, and both the physical batch and the corresponding
closed source interval lie in the anchor window $[-s,2s-1]$.
At the terminal plateau the preceding source step adds at most one unit.
These are consequences of the actual simulation invariant in
\cite{HighRateStorage2026}, not of an idealized rescaled schedule.

\begin{lemma}[Guarded path projection]\label{lem:guarded-projection}
\formal{TransferNativeProjection.lean}{TransferGuards.native_guarded_projection}
For every source vertex set $B$, there is a simulator set $E$ with
\[
 |E|\le [4B_0+k(6s+1)]|B|
\]
such that every surviving simulator dependency path projects to a source
path with endpoints $q(u),q(v)$ avoiding $B$, endpoints included.
Moreover $|y_h(u)-x_h(q(u))|\le3s$ for every corresponding head, including
the input head.
\end{lemma}
\begin{outline}
The checked execution invariant of the actual compiler proves the position
bound, including the terminal batch and drain. Any two time-ordered source
visits in a closed band admit a backward dependency path within that band,
even if the source walk leaves both sides between them. Follow the state
edge when its predecessor stays in the band; otherwise jump to the actual
last prior visit to the current boundary. Unit speed supplies that visit.
Time decreases at each step. Thus endpoint shadow alone projects a simulator
work hop at cell $z$ to a source path in $[z-3s,z+3s]$. Equal projected
endpoints require only a singleton path.

Mark every whole physical batch whose closed projected source interval meets
$B$. A source cut time belongs to at most four such batches, including the
extra terminal batch and drain. Each costs at most $B_0$. Also mark the later
endpoint of a simulator work hop whenever its projected interval strictly
straddles a cut time $b$ and its cell is within $3s$ of $x_h(b)$. For each
$(b,h,z)$ at most one consecutive-visit interval straddles $b$: such intervals
have disjoint interiors and $q$ is nondecreasing. This costs at most
$k(6s+1)|B|$. These are the sets whose union is $E$.

A surviving state edge projects either to a singleton or to a source spine
interval disjoint from $B$. A surviving work edge has the band path just
constructed. A cut vertex on that path would either be a projected endpoint,
contradicting the whole-batch guard, or lie strictly between them, contradicting
the hop guard. Concatenation proves the path statement. This uses every actual
work last-visit edge and does not introduce input-revisit edges.
\end{outline}

Let $\alpha_i,u_i$ be the actual source and simulator append times for output
ordinal $i$. Ordinals are zero-based and matched even when output bit values
repeat. Let $O(v)$ count source appends strictly before time $v$.

\begin{lemma}[Output look-ahead and cones]\label{lem:cone-transport}
\formal{TransferNativeCones.lean}{TransferCones.native_cone_transport_interior,TransferCones.native_cone_transport}
For every existing ordinal, $\alpha_i\le q(u_i)$ and
$O(q(u_i)+1)\le i+2s$. Put
\[
 G_s=4B_0+k(6s+1)+2s+1.
\]
For every $B$ there is a guard $\Gamma$ of size at most $G_s|B|+2s$.
Every append root $u_i\notin\Gamma$ has $i+2s$ defined, and its root-exempt
input cone satisfies
\[
 \operatorname{Cone}_{N_s}(\Gamma,i)
 \subseteq \operatorname{Dil}_{3s}\bigl(\operatorname{Cone}_M(B,i+2s)\bigr)
          \cup\operatorname{Boundary}_{3s}(P).
\]
The dilation is in valid raw coordinates; the boundary comprises the first
and last $3s$ positions, with overlap allowed.
\end{lemma}
\begin{outline}
At physical batch start $j$, the emitted output count is
$E_j=s\lfloor O(js)/s\rfloor$ and fewer than $s$ source outputs are pending.
The compiler emits a prefix of length either zero or $s$ and retains the
suffix. Thus its FIFO order gives $\alpha_i\le q(u_i)$. Through the source
step $q(u_i)$, inclusive, at most $s+1$ new outputs follow $O(js)$.
For a batch emission $i\ge E_j$, this gives
$O(q(u_i)+1)\le E_j+(s-1)+(s+1)\le i+2s$.
The drain has fewer than $s$ residual outputs and satisfies the same bound.
These actual append-order facts are checked in
\lean{TransferFIFO.actual_sourceRoot_le_clock} and
\lean{TransferFIFO.actual_ordinal_lookahead}.
Hence $q(u_i)<\alpha_{i+2s}$ when that future ordinal exists.

Take the guard of \cref{lem:guarded-projection}. Add the last at most $2s$
append roots and every root $u_i$ whose closed interval
$[\alpha_i,\alpha_{i+2s}]$ meets $B$. A cut time belongs to at most $2s+1$
these intervals: their smallest and largest indices differ by at most $2s$,
since append times increase strictly. This proves the guard bill. For a
surviving root, prepend the $B$-free source spine from $\alpha_{i+2s}$ to
$q(u_i)$ to the projected path. At a raw endpoint, the input-head shadow
places the simulator raw address within $3s$ of the source address. If that
source address is valid it belongs to the source cone. Otherwise the original
valid simulator address is in the displayed boundary. Singleton paths and
endpoints in cuts are included. The look-ahead bounds output counts, not the
physical waiting time of a pending output bit.
The formalization also gives the version before adding the terminal roots:
the guard costs at most $G_s|B|$, and the inclusion holds whenever the
future ordinal exists. This is the version used for interior entry lanes.
Every constructed guard consists of actual simulator vertices.
\end{outline}

\subsection{Finite raw-input cuts and erosion}

For $U\subseteq[0,P)\cap\Z$, regard $U$ as a subset of the entire integer
line and define
\[
 E_s(U)=\{z\in\Z:[z-3s,z+3s]\cap\Z\subseteq U\},\qquad
 C_s=4B_0+k(6s+1)+2s-1.
\]
Thus erosion pays the physical input boundaries as well as internal gaps.

\begin{lemma}[Finite interval erosion]\label{lem:input-erosion}
\formal{InputErosion.lean}{interval_union_erosion_loss,demanded_erosion_loss}
If a finite integer set $U$ is a union of $r$ integer intervals, then
$|U\setminus E_s(U)|\le6sr$. Consequently, for every $V\subseteq U$,
$|V|\le |V\cap E_s(U)|+6sr$. The intervals need not be disjoint.
\end{lemma}
\begin{outline}
Every point of an interval outside its two end strips of width $3s$ has
its whole radius-$3s$ neighborhood in that interval, hence in $U$.
Therefore $U\setminus E_s(U)$ lies in the union of the $2r$ end strips,
whose total cardinality is at most $6sr$. The second inequality partitions
$V$ into its eroded and noneroded parts. The Lean statement allows an
arbitrary nonnegative integer radius and a finite interval cover.
\end{outline}

\begin{proposition}[Finite eroded-input transfer]\label{prop:eroded-input-transfer}
\formal{TransferErodedInvoice.lean}{TransferEroded.native_eroded_fibre_invoice,TransferEroded.native_eroded_fibre_invoice_cover}
On an actual source run and its aligned simulator, a source cut $B$ from
any set of existing append ordinals to all source reads of $U$ lifts to
a simulator cut $\Gamma$ from the matched append roots to all simulator
reads of $E_s(U)$, with
\[
 |\Gamma|\le C_s|B|+2s-1.
\]
This finite geometric assertion needs no correctness or common source trace.
On a nonempty common simulator trace fibre $S$ of correctly transposed inputs,
put $D_S=P-\log_2|S|$ and let $v_s$ be the simulator's complete
step-value width, including its output field. For every $X\in S$,
every $V\subseteq U\subseteq\{0,\ldots,P-1\}$ and every original cut $B$ separating the demanded
roots for $V$ from all original reads of $U$,
\[
|V\cap E_s(U)|\le D_S+v_s(C_s|B|+2s-1).
\]
If $U$ is a union of $r$ integer intervals, then
$|U\setminus E_s(U)|\le6sr$; in particular the left side can be replaced
by $|V|-6sr$. These finite assertions hold at every positive entry width.
\end{proposition}
\begin{outline}
The physical input head satisfies the same shadow estimate
$|y_0(u)-x_0(q(u))|\le3s$ used in \cref{lem:guarded-projection}.
In an ordinary core both positions lie in the aligned anchor plus
$[-s,2s-1]$; the terminal replacement by the last live source step costs
at most one unit, and the stationary drain has error at most $s$.
These estimates apply to the actual input-head coordinate, including the
terminal core. Therefore a simulator read at a point of $E_s(U)$ projects
to an actual source read in $U$. Its connector is the singleton source
vertex $q(u)$; no input-revisit edge is introduced.

Use the edge guards of \cref{lem:guarded-projection}, costing at most
$[4B_0+k(6s+1)]|B|$. In the source graph with $B$ deleted, let $R_U$
be the vertices reaching a read of $U$. Its intersection with a $B$-free
state-spine interval is an upper terminal segment, so there are at most
$|B|+1$ positive activations. At an activation time $c$, the FIFO invariant
in \cref{lem:cone-transport} leaves at most $2s-1$ output ordinals whose
actual source root precedes $c$ but whose simulator shadow is at least $c$.
Delete those matched simulator roots. This costs
$(2s-1)(|B|+1)$. Every surviving simulator path to an eroded read projects
to a source path avoiding $B$ and ending at a read of $U$. Its projected
root lies in $R_U$, whereas its actual source root lies outside $R_U$ by
separation; the activation guard forbids this. Singleton paths and cut
roots are covered by the endpoint-inclusive guards. If the source has no
successful step there are no append ordinals and the assertion is vacuous.
The complete geometric lift, including this zero-time case, is
\lean{TransferEroded.native_eroded_cut_lift}.

Restrict the demanded source coordinates to $V\cap E_s(U)$ and apply
the lift. Its resulting separator is on the one simulator trace shared
by $S$. The frozen-trace injection of \cref{prop:frozen-invoice} gives
$|V\cap E_s(U)|\le D_S+v_s|\Gamma|$, proving the displayed bound.
The weaker coefficient $v_s+1$ is an immediate consequence; a full step
record already includes the output field and needs no further root surcharge.
This direct injection also explains the all-width finite
claim; no canonical-width asymptotic theorem is used. The cut can depend
on $X,U,V$ because the frozen-trace assertion holds for every cut on that
fixed fibre, with no restriction to inputs sharing the original trace.

Apply \cref{lem:input-erosion} to the given interval cover. The maximal
interval components are one such cover.
Highly fragmented sets may have empty erosion; this cost is not discarded.
\end{outline}

\subsection{An original entry quotient carries a large storage code}

For one reference source run $X$ and cut $B$, let $Q_X^B(e)$ contain entry $f$
if any bit of $f$ belongs to the source root-exempt cone of any of the $b$
output bits of entry $e$. This is the full quotient: all lanes and self-edges
are retained. For a relation $Q$ on $J$, let $\operatorname{Cap}_m(Q)$ be the
base-two logarithm of the largest storage-code cardinality over $m$-bit symbols.
Recovery means that agreement on all helpers in $Q(e)$ forces agreement at $e$.
Write $\Delta_m(Q)=m|J|-\operatorname{Cap}_m(Q)$; both are zero for $m=0$.

\begin{theorem}[Theorem Q: capacity of original entry relations]\label{thm:original-entry-capacity}
\formal{TransferEntryNative.lean}{TransferEntry.native_original_entry_capacity}
Let $S$ be a nonempty set of correctly transposed inputs with one common complete
simulator trace. Set $m=(b-6s)_+$, $D_S=P-\log_2|S|$, and let $v_s$ be the
complete step-value width of the simulator. For every reference $X\in S$,
every source vertex set $B$ on that run and every entry set $J$,
\[
 \Delta_m(Q_X^B[J])\le D_S+v_sG_s|B|.
\]
The extracted code need not contain $X$, and its members need not have the same
original trace. The finite statement holds at every $b\ge1$.
\end{theorem}
\begin{outline}
If $m=0$ the conclusion is immediate. Otherwise use the guard in
\cref{lem:cone-transport} without the final-output tail, of size at most
$G_s|B|$. An interior lane $3s\le\gamma<b-3s$ is never in that tail. Partition
$S$ by the joint complete guard values and every raw bit outside the interior
lanes of $J$. There are at most $2^{v_sG_s|B|+P-m|J|}$ labels. A largest cell,
projected to the remaining lanes, is an injective image and has logarithmic
cardinality at least $m|J|-D_S-v_sG_s|B|$.

For a guarded interior root its output bit is fixed by the complete guard
value. For another root, the output look-ahead stays within entry $e$. If an
interior raw helper survives the fixed exterior, its distance-$3s$ source
neighbor stays within the same helper entry. Thus every nonconstant helper
is in $Q_X^B(e)\cap J$. The fixed-trace cone-determination theorem of
\cite{MacroscopicCrossings2026} gives recovery of each bit throughout this
cell. The checked cone-determination step is
\lean{nativeInputCone_determines}; guarded roots use the output field already
present in the complete record. Boundary helpers are fixed. The reference run fixes $B,Q_X^B$ and the
simulator guard, whereas only the simulator trace is required to be common.
This proves the storage-code property and the capacity inequality.
\end{outline}

\begin{corollary}[One event for all original capacity certificates]\label{cor:entry-success-tail}
\formal{TransferEntryNative.lean}{TransferEntry.uniform_entry_success_tail,TransferEntry.certified_original_cut_lower_bound}
Fix a machine, a linear time constant, canonical width $b(d)$, and
$0<\delta<1$. Some fixed $s$ and all sufficiently large $d$ have the following
property: fewer than $2^{(1-\delta/2)P}$ timely correct inputs violate any of
\[
 \Delta_{b-6s}(Q_X^B[J])\le\delta P+v_sG_s|B|,
 \qquad\text{for all }B,J.
\]
If $0<\alpha,\rho\le1$, $|J|\ge\alpha d^2$ and
$\operatorname{Cap}_{b-6s}(Q_X^B[J])\le(1-\rho)(b-6s)|J|$, choosing
$\delta=\rho\alpha/4$ gives, on the good event at large $d$,
$|B|\ge\rho b|J|/(2v_sG_s)$.
\end{corollary}
\begin{outline}
The successful-subset simulator resource theorem of~\cite{HighRateStorage2026}
bounds the number of attained simulator traces by $2^{aP/s}$ for a fixed
$a$. Choose $s$ with $a/s\le\delta/2$ before dimension grows. Fibres of size
at least $2^{(1-\delta)P}$ satisfy every finite inequality above. The total
number of successful inputs in smaller attained fibres is strictly below
$2^{(1-\delta/2)P}$. There is no union bound over $B,J$ or reference inputs.
For the final assertion take $b\ge24s$: the deficit is at least
$3\rho b|J|/4$, whereas $\delta P\le\rho b|J|/4$, and rearrange.
\end{outline}

Acyclic entry relations have singleton storage codes and are a special case.
Bit-level acyclicity is insufficient: shifting a target by $2s$ can introduce
self-dependence within an entry. No separation, degree or missing-self-edge
hypothesis is supplied automatically. A capacity theorem is useful here only
after its own premises have been established on the actual quotient.

\subsection{Expansion from every large output set to original crossings}

Let $\mathrm{CL}_h$ be the earlier endpoints of same-tape crossing consecutive
revisits whose two excursions exceed $2h$, with the strict crossing order of
\cite{MacroscopicCrossings2026}.

\begin{theorem}[Theorem E: original adaptive crossing expansion]\label{thm:original-crossing-expansion}
\formal{TransferCrossingTail.lean}{TransferCrossing.uniform_original_crossing_expansion}
Fix a native machine, $C\ge0$ and $\alpha>0$.
At canonical width there are $\eta,\lambda,\beta>0$ such that, at every
sufficiently large $d$, fewer than $2^{(1-\eta)P}$ timely correct inputs fail
the following simultaneous statement about their original execution:
for every $|V|\ge\alpha P$, every vertex cut from the actual append roots
$\operatorname{Out}_M(V)$ to $\mathrm{CL}_{\lfloor\lambda P\rfloor}$ has
size at least $\beta|V|$. No original-trace heaviness assumption is imposed.
\end{theorem}
\begin{outline}
We first give the finite crossing-target addition to
\cref{lem:guarded-projection}. Let $H\ge h\ge1$, and let $K_0$ bound the
number of grid-$H$ epochs of each simulator head. Each epoch has diameter
less than $2H$. A crossing pair $p<p'<u<u'$ has opposite polarities after
one choice of sign: its first hop is above base $a$, and its second below
base $b>a$. Choose an epoch containing a maximum of the first hop and one
containing a minimum of the second. Their fixed signed extremizers $t_1,t_2$
satisfy $p<t_1\le p'<u\le t_2<u'$. Let $t_0$ minimize the signed walk before
$t_1$, and $t_3$ maximize it after $t_2$. The resulting zigzag has overlap
\[
 [L_0,U_0]=[\max(y(t_0),y(t_2)),\min(y(t_1),y(t_3))]
\]
containing both bases. For fixed sign, ordered epoch pair and endpoint role,
each base level determines at most one endpoint: distinct consecutive-visit
hop interiors at one level cannot both contain the fixed extremizer.

Put $R=3s$ and $w=2h+2R+2$. Discard targets whose assigned level is within
$w$ of either end of the overlap. The preceding injection bounds their number
by $4kK_0^2(w+1)$. Indeed the two extremizer epochs are distinct: their
values differ by more than $2H$, whereas one epoch has diameter less than
$2H$. The checked label records the sign, endpoint role, two epoch indices
and a level offset at most $w$. Reversing the strictly ordered epoch pair
encodes the upper end strip, so no extra strip bit is needed. The label
space has size $4K_0^2(w+1)$ per tape, and actual consecutive-visit
uniqueness makes the encoding injective.
At a remaining first-role target of level $a$, take
$\ell=a-R-1$ and $c=\ell+2h+1$. The shadowed zigzag visits below $\ell$,
above $c$, below $\ell$, then above $c$. The last $\ell$ visit before
$q(v)$ and first after it form a genuine upper hop; the neighboring $c$ visits
supply a crossing lower hop. Both excursions exceed $2h$.
For a second-role target at $b$, reflect the construction with
$\ell=b+R+1$, $c=\ell-2h-1$; the crossing partner occurs \emph{before} the
chosen lower hop. In either case the checked band-path argument connects
$q(v)$ backward to the appropriate source endpoint within the common band
$[y(v)-R-1,y(v)+R+1]$ in the unsigned coordinates.

Guard a good target whenever its connector band contains a cut position.
At a fixed cell, a hop of excursion $>2H$ uses at least $4H+2$ transitions;
their interiors are disjoint. This guard therefore costs at most
$k(6s+3)\lfloor(U-1)/(4H+2)\rfloor|B|$.

There is one further root bill. Let $R_B$ be the source vertices reaching
$\mathrm{CL}_h$ while avoiding $B$. On each $B$-free spine interval it is an
upper terminal segment, so it has at most $|B|+1$ positive activations.
At an activation $c$, a delayed ordinal satisfies
$\alpha_i<c\le q(u_i)$. Thus $i<O(c)$, while the checked pre-step FIFO
bound gives $O(c)\le O(q(u_i))\le i+2s-1$.
There are at most $2s-1$ such ordinals. Delete their matched simulator roots.
Any remaining simulator path to a surviving crossing target projects to a
source path avoiding $B$ and continues along its connector. Its projected
root is in $R_B$, whereas its separated actual source root is outside $R_B$;
the activation guard forbids this mismatch. Thus a source separator lifts to
one of size at most $A|B|+F$, where
\[
\begin{split}
 A&=4B_0+k(6s+1)+k(6s+3)\lfloor(U-1)/(4H+2)\rfloor+2s-1,\\
 F&=4kK_0^2(2h+6s+3)+2s-1.
\end{split}
\]
This exact finite lift is \lean{TransferCrossing.native_crossing_cut_lift}.

Finally use the successful-subset trace bound of~\cite{HighRateStorage2026}
and the correct-subset expansion theorem of~\cite{MacroscopicCrossings2026}.
Its permitted entropy budget depends only on tape count, the simulator time
coefficient and $\alpha$, so it can be chosen before the fixed block size.
On every
remaining heavy successful fibre, the latter theorem gives a fixed
$H=\lfloor\theta P\rfloor$ and cut lower bound $e|V|$ simultaneously for
all $|V|\ge\alpha P$, as checked in
\lean{TransferCrossing.native_crossing_fibre_bound}.
Since $U=O(P)$ and $\theta P\le2H$ at large $P$, one fixed $K_0$ bounds
the epoch counts and one fixed $\overline A$ bounds $A$.
Choose $h=\lfloor\lambda P\rfloor$ with $0<\lambda\le\theta$ sufficiently small that
$F\le e\alpha P/2$ at large $d$. Then
$e|V|\le A|B|+F$ gives $|B|\ge e|V|/(2\overline A)$.
The simultaneous choices and bounds are proved in
\lean{TransferCrossing.exists_uniform_crossing_budgets}.
The simulator has at most $2^{\eta P}$ attained successful traces after
choosing its fixed block size, and the fibres below the required
$2^{(1-2\eta)P}$ threshold together contain fewer than
$2^{(1-\eta)P}$ inputs. Every cut and output set is covered on the same
remaining event. This proves the asserted quantifiers and the ambient bound.
\end{outline}

Theorem E concerns arbitrary macroscopic output sets at canonical widths;
the direct off-carrier theorem concerns its selected slab demands at all
widths. Neither subsumes the other. Theorem Q adds a different certificate:
a cheap cut would need an original entry relation with almost full storage
capacity. Finding a cut and a nonvacuous deficit together remains open.
